% Memoria de la Práctica 1 (PLN I, 2026-27). Formato ACL, máximo 4 páginas sin referencias.
%
% Compilación (desde esta carpeta):   latexmk -pdf main.tex
% Requiere acl.sty y acl_natbib.bst del repositorio oficial acl-org/acl-style-files.
%
% Convención de marcadores (todo lo que salga en rojo debe desaparecer antes de entregar):
%   \pend{texto}  -> dato o cifra pendiente de confirmar o de obtener
%   \todo{texto}  -> párrafo por redactar cuando haya resultados
% Comprobación previa a la entrega:  grep -rn "pend{\|todo{" sections main.tex
\documentclass[11pt]{article}

\usepackage[preprint]{acl}
\usepackage{times}
\usepackage{latexsym}
\usepackage[T1]{fontenc}
\usepackage[utf8]{inputenc}
\usepackage{microtype}
\usepackage{inconsolata}
\usepackage[spanish,es-tabla,es-nodecimaldot]{babel}

\usepackage{amsmath,amssymb}
\usepackage{booktabs}
\usepackage{graphicx}
\usepackage{xcolor}

% ---------- Macros de trabajo ----------
\newcommand{\pend}[1]{\textcolor{red}{#1}}
\newcommand{\todo}[1]{\textcolor{red}{[\textbf{PENDIENTE}: #1]}}
% Inserta la figura si existe en figures/; si no, deja un recuadro visible.
\newcommand{\figorph}[3]{%
  \IfFileExists{figures/#1}%
    {\includegraphics[width=#2]{figures/#1}}%
    {\fbox{\parbox{#2}{\centering\small\textcolor{red}{Figura pendiente:\\#3}}}}}

% Análisis interno (logit lens, atención): activar con \internotrue cuando haya puntos de control.
\newif\ifinterno
\internotrue

% ---------- Notación ----------
\newcommand{\hid}[1]{\mathbf{h}_{#1}}
\newcommand{\LN}{\mathrm{LN}}
\newcommand{\MHA}{\mathrm{MHA}}
\newcommand{\FFN}{\mathrm{FFN}}

\title{Profundidad por recurrencia: un \emph{Looped Transformer} frente a GPT\\
en problemas aritméticos en castellano de dificultad controlada}

\author{\textbf{Alejandro García Belmonte}\\Universitat d'Alacant
\And \textbf{Stanislav Gatin}\\Universitat d'Alacant
\AND \textbf{Rodrigo Gavilán Salmerón}\\Universitat d'Alacant
\And \textbf{Pau Mateo Lillo}\\Universitat d'Alacant}

\begin{document}
\maketitle

\begin{abstract}
Un \emph{Looped Transformer} aplica $N$ veces un único bloque con los mismos
pesos y obtiene profundidad de cómputo sin parámetros adicionales. Lo
implementamos desde cero junto a un GPT, con una interfaz común validada
numéricamente, y los entrenamos en problemas aritméticos en castellano de 1 a 6
operaciones, reservando los niveles 5 y 6 para medir la generalización.
Comparamos ambas familias a igual profundidad, parámetros y cómputo, y
estudiamos el interior del bucle con \emph{logit lens} y métricas de
representación y atención. Con ocho iteraciones, un Looped Transformer de
0.97\,M de parámetros iguala o supera a un GPT de ocho capas con 6.5\,M, sobre
todo en problemas con más operaciones que los vistos al entrenar, y las
iteraciones necesarias crecen con la dificultad; con dos y cuatro queda al nivel
de un GPT de igual profundidad o algo por debajo. Código:
\url{https://github.com/rodrigogavsal/looped-math}.
\end{abstract}

\section{Introducción}

Los Transformers resuelven los problemas de varios pasos aplicando una sucesión
de capas, cada una con sus propios parámetros. La profundidad parece clave para
el razonamiento \citep{saunshi2025reasoning}, pero añadir capas encarece de forma
lineal los parámetros, la memoria y el entrenamiento. El Universal
Transformer \citep{dehghani2019universal} ya aplica de forma recurrente un mismo
bloque; trabajos recientes lo estudian como \emph{Looped Transformer}: un único
bloque aplicado $N$ veces sobre su propia salida, con profundidad de cómputo
creciente y parámetros constantes
\citep{giannou2023looped,saunshi2025reasoning,geiping2025scaling}.

Estudiamos esta propuesta a pequeña escala y con control total de los datos:
entrenamos desde cero GPT y Looped Transformers de pocos millones de parámetros
en problemas aritméticos en castellano cuya dificultad es el número de
operaciones. Buscamos responder a \textbf{(P1)}~si un bucle de $N$ iteraciones
iguala a un GPT de $N$ capas con hasta $N$ veces menos parámetros (al compartir el bloque),
\textbf{(P2)}~si las iteraciones necesarias crecen con la dificultad y
\textbf{(P3)}~si el modelo generaliza a problemas con más operaciones que los
vistos al entrenar.\ifinterno{} Además, \textbf{(P4)}~analizamos el interior del bucle con \emph{logit lens} y métricas de representación y atención.\fi{}

Aportamos una implementación propia, sobre el código de \citet{raschka2024build},
de ambos modelos con una interfaz común y validada numéricamente
(Sec.~\ref{sec:validacion}); un conjunto de problemas sintéticos con seis niveles
y los dos más difíciles reservados para evaluación (Sec.~\ref{sec:datos}); y una
comparación controlada a igual cómputo y a igual número de parámetros, con
barrido de iteraciones y análisis interno del bucle
(Secs.~\ref{sec:experimentos} y~\ref{sec:analisis}).

\section{Trabajo relacionado}

\paragraph{Pesos compartidos y bucles.}
El Universal Transformer \citep{dehghani2019universal} aplica de forma
recurrente un mismo bloque sobre la secuencia. \citet{giannou2023looped}
muestran que un Transformer con bucles puede ejecutar algoritmos como un
ordenador programable, y \citet{saunshi2025reasoning} comparan, a igual
profundidad efectiva, modelos con y sin bucles en tareas sintéticas de
razonamiento. A mayor escala, \citet{geiping2025scaling} escalan el cómputo en
inferencia iterando un bloque latente. Nos distinguimos por el dominio, el barrido de iteraciones y el análisis interno.


\section{Método}

\subsection{Arquitecturas}
\label{sec:arquitecturas}

Ambos modelos comparten el bloque de GPT-2 \citep{radford2019language}:
normalización previa, atención multicabeza causal, red \emph{feed-forward} de
dimensión $4d$ con GELU y conexiones residuales,
$B_\theta(\hid{}) = \mathbf{u} + \FFN(\LN_2(\mathbf{u}))$ con
$\mathbf{u} = \hid{} + \MHA(\LN_1(\hid{}))$. La entrada es la suma de los
\emph{embeddings} de carácter y de posición (aprendido), $\hid{0}$, y la salida
$\mathrm{softmax}(W_o\,\LN(\hid{D}))$ con $D$ la profundidad efectiva. Un
\textbf{GPT} de $L$ capas usa bloques distintos, $\hid{\ell}=B_{\theta_\ell}(\hid{\ell-1})$;
un \textbf{Looped Transformer} de $N$ iteraciones reutiliza el mismo,
\begin{equation}
\hid{t} = B_\theta(\hid{t-1}),\qquad t = 1,\dots,N,
\end{equation}
con $D=L$ y $D=N$ respectivamente. Con dimensión $d$, vocabulario $V$ y contexto
$T$, un bloque tiene $P_B = 12d^2 + 10d$ parámetros y el resto del modelo
$P_E = (2V+T)d + 2d$, de modo que $P_\text{GPT}(L) = L\,P_B + P_E$ y
$P_\text{Looped}(N) = P_B + P_E$ no depende de $N$, mientras que el cómputo por
token, $F(D) \approx D\,(24d^2 + 4Td) + 2dV$, crece con $D$ en ambos.

\subsection{Implementación y validación}
\label{sec:validacion}

Partimos del código de \citet{raschka2024build}, con la atención explícita para
devolver sus pesos; ambos modelos exponen la misma interfaz (logits, estados
ocultos y atenciones por paso). Comprobamos la corrección con pruebas
automáticas: Looped con $N{=}1$ coincide con GPT con $L{=}1$ y mismos
pesos ($10^{-6}$); Looped con $N\in\{2,4,8,12\}$ coincide, en logits y en todos
los estados ocultos, con un GPT de $N$ capas con copias del bloque compartido
($10^{-5}$); el número de parámetros no depende de $N$ y la atención es causal.

\subsection{Datos}
\label{sec:datos}

Cada ejemplo es una historia breve en castellano cuyo nivel es el número de
operaciones, de 1 a 6: suma, resta, doble, triple, mitad y tercera y cuarta
parte, con números de hasta 20 en el enunciado y respuestas entre 1 y 99. Hay
relatos con un protagonista o un recipiente (tipo~A) y problemas de comparación
entre varias personas (tipo~B) cuyas frases pueden estar desordenadas, en
estilo detallado o compacto y a partir de plantillas. Un ejemplo de nivel~2 con frases desordenadas:
\begin{quote}\small
\emph{Julia tiene la mitad que Laura. Laura tiene el triple de postales que Pablo.
Pablo tiene 18 postales. Entonces, ¿cuántas postales tiene Julia?}\ \ $\to$ 27
\end{quote}
Cada ejemplo guarda su traza de operaciones, con la que se verifica la
respuesta, y ninguna pregunta de validación o de test aparece en el
entrenamiento. Los niveles 1--4 se reparten en entrenamiento (200\,000
ejemplos, 50\,000 por nivel), validación (4\,000) y test en distribución
(ID, 8\,000, 2\,000 por nivel). Un segundo test de plantillas nuevas (8\,000
ejemplos) usa 18 redacciones reservadas que no aparecen en el entrenamiento.
Los niveles 5 y 6 (4\,000 ejemplos) solo se usan como test fuera de
distribución (OOD). Las heurísticas sin razonamiento (respuesta más frecuente, primer o último número) aciertan entre el 0.6 y el 7.6\,\%.

\subsection{Entrenamiento y evaluación}
\label{sec:entrenamiento}

El texto de entrada es \texttt{Pregunta: \dots\textbackslash nRespuesta: } seguido
del número y de un fin de secuencia, tokenizado por caracteres (97 símbolos); la
pérdida se calcula solo sobre la respuesta. Usamos AdamW
\citep{loshchilov2019decoupled} ($\beta=(0.9, 0.95)$, decaimiento 0.1), tasa
$10^{-3}$ con calentamiento de 200 pasos y decaimiento coseno hasta el 10\,\%,
lotes de 64, recorte de gradiente 1.0 y precisión mixta, durante 42\,000 pasos y con tres semillas (0, 1 y 2).
Evaluamos por generación voraz (hasta 6 caracteres, primer entero extraído) y
medimos la exactitud por nivel (coincidencia exacta); en los modelos con bucles variamos las
iteraciones en test, $N'$, de 1 a $2N$ sin reentrenar.

\section{Experimentos}
\label{sec:experimentos}

\begin{table*}[!t]
\centering\small
\setlength{\tabcolsep}{4pt}
\begin{tabular}{lrccccccc}
\toprule
 & & \multicolumn{4}{c}{En distribución} & \multicolumn{2}{c}{OOD} & \\
\cmidrule(lr){3-6}\cmidrule(lr){7-8}
Modelo & Parám. (M) & N1 & N2 & N3 & N4 & N5 & N6 & Plant. nuevas \\
\midrule
GPT $L{=}1$ & 0.97 & $87.6{\scriptstyle\pm0.6}$ & $35.3{\scriptstyle\pm0.3}$ & $8.9{\scriptstyle\pm0.6}$ & $6.7{\scriptstyle\pm0.6}$ & $5.1{\scriptstyle\pm0.4}$ & $4.1{\scriptstyle\pm0.3}$ & $16.1{\scriptstyle\pm0.8}$ \\
GPT $L{=}2$ & 1.76 & $96.8{\scriptstyle\pm0.5}$ & $73.1{\scriptstyle\pm3.5}$ & $26.1{\scriptstyle\pm3.4}$ & $16.9{\scriptstyle\pm1.7}$ & $12.5{\scriptstyle\pm0.9}$ & $8.9{\scriptstyle\pm0.4}$ & $22.4{\scriptstyle\pm2.4}$ \\
GPT $L{=}4$ & 3.34 & $96.9{\scriptstyle\pm0.2}$ & $78.7{\scriptstyle\pm1.8}$ & $48.8{\scriptstyle\pm4.1}$ & $34.9{\scriptstyle\pm4.1}$ & $23.9{\scriptstyle\pm3.0}$ & $14.5{\scriptstyle\pm1.8}$ & $25.5{\scriptstyle\pm1.9}$ \\
GPT $L{=}8$ & 6.49 & $97.0{\scriptstyle\pm0.2}$ & $\mathbf{80.9}{\scriptstyle\pm0.4}$ & $54.3{\scriptstyle\pm2.1}$ & $40.2{\scriptstyle\pm1.8}$ & $25.4{\scriptstyle\pm2.7}$ & $13.6{\scriptstyle\pm1.7}$ & $26.7{\scriptstyle\pm0.6}$ \\
\midrule
Looped $N{=}2$ & 0.97 & $96.9{\scriptstyle\pm0.4}$ & $71.9{\scriptstyle\pm0.4}$ & $25.6{\scriptstyle\pm1.2}$ & $16.5{\scriptstyle\pm0.2}$ & $12.2{\scriptstyle\pm0.4}$ & $9.0{\scriptstyle\pm0.4}$ & $21.1{\scriptstyle\pm1.7}$ \\
Looped $N{=}4$ & 0.97 & $\mathbf{97.3}{\scriptstyle\pm0.2}$ & $78.7{\scriptstyle\pm0.6}$ & $47.0{\scriptstyle\pm1.1}$ & $31.6{\scriptstyle\pm0.3}$ & $21.1{\scriptstyle\pm2.3}$ & $14.0{\scriptstyle\pm0.1}$ & $23.3{\scriptstyle\pm1.0}$ \\
Looped $N{=}8$ & 0.97 & $97.1{\scriptstyle\pm0.2}$ & $79.0{\scriptstyle\pm0.2}$ & $\mathbf{57.2}{\scriptstyle\pm1.1}$ & $\mathbf{45.8}{\scriptstyle\pm1.4}$ & $\mathbf{33.4}{\scriptstyle\pm2.2}$ & $\mathbf{21.2}{\scriptstyle\pm2.2}$ & $\mathbf{29.8}{\scriptstyle\pm4.2}$ \\
\midrule
GPT $L{=}4$ ($d{=}128$) & 0.88 & $96.5{\scriptstyle\pm0.5}$ & $65.5{\scriptstyle\pm7.4}$ & $35.0{\scriptstyle\pm10.7}$ & $24.3{\scriptstyle\pm8.7}$ & $16.5{\scriptstyle\pm5.1}$ & $11.3{\scriptstyle\pm3.0}$ & $22.6{\scriptstyle\pm2.5}$ \\
\bottomrule
\end{tabular}

\caption{Exactitud (\%) por nivel, media $\pm$ desviación típica sobre tres
semillas (niveles 1--4: test en distribución; 5 y 6: fuera de distribución;
\emph{Plant. nuevas}: redacciones no vistas). En negrita, la mejor media.}
\label{tab:main}
\end{table*}

\begin{figure*}[!t]
\centering
\figorph{loops_sweep.pdf}{\textwidth}{exactitud frente al número de iteraciones usado en test}
\caption{Exactitud de un Looped Transformer entrenado con $N$ iteraciones al
evaluarlo con $N'$ sin reentrenar (banda: mínimo y máximo entre semillas).}
\label{fig:sweep}
\end{figure*}

\begin{figure*}[!t]
\centering
\figorph{internal_main.pdf}{\textwidth}{logit lens, atención a los dígitos y dinámica de representación y atención}
\caption{Interior de Looped $N{=}8$ y GPT $L{=}8$ con \emph{teacher forcing}.
Arriba: fracción de problemas con la respuesta completa correcta al proyectar
cada estado con la cabeza de salida (\emph{logit lens}) y atención de la posición
de respuesta a los dígitos del enunciado, por nivel. Abajo: cambio relativo,
rango efectivo, entropía normalizada y divergencia JS entre pasos (niveles 1--4;
banda: mínimo y máximo entre semillas).}
\label{fig:interno}
\end{figure*}

\subsection{Configuraciones y comparaciones controladas}

Todos los modelos usan contexto de 512 caracteres, vocabulario de 97 símbolos y
$d{=}256$ con 8 cabezas, salvo un GPT de 4 capas con $d{=}128$ y 4 cabezas
(contraste C3). Entrenamos GPT con
$L\in\{1,2,4,8\}$ capas y Looped Transformers con $N\in\{2,4,8\}$ iteraciones,
con las semillas 0, 1 y 2 (24 entrenamientos de 42\,000 pasos).
La Tabla~\ref{tab:main} indica el número de parámetros de cada modelo. La matriz
permite tres contrastes: \textbf{(C1)}~\emph{igual cómputo}, Looped $N$ frente a GPT
$L{=}N$, con los mismos FLOPs por token y $N$ veces más parámetros el segundo;
\textbf{(C2)}~\emph{iguales parámetros}, Looped $N$ frente a GPT $L{=}1$, que no
comparten el cómputo; y \textbf{(C3)}~\emph{parámetros y profundidad parecidos},
Looped $N{=}4$ (0.97\,M de parámetros) frente a ese GPT de 4 capas con $d{=}128$
(0.88\,M), que tiene la misma profundidad y casi los mismos parámetros pero 3.2
veces menos FLOPs por token.

\subsection{Resultados}

En las comparaciones, la primera cifra es siempre la del Looped y la segunda la
del GPT.

\paragraph{Igual número de parámetros (C2).}
Con 0.97\,M de parámetros, el Looped Transformer llega a 25.6 y 16.5\,\% en los
niveles 3 y 4 con $N{=}2$, a 47.0 y 31.6\,\% con $N{=}4$ y a 57.2 y 45.8\,\% con
$N{=}8$ (Tabla~\ref{tab:main}), mientras que el GPT de una capa apenas resuelve
esos niveles (8.9 y 6.7\,\%, frente al 6.2 y 5.4\,\% de responder siempre la
respuesta más frecuente): iterar el bloque aporta una capacidad que no se
obtiene con más parámetros.

\paragraph{Igual cómputo (C1).}
Con $N{=}2$ y $N{=}4$ el Looped queda al nivel del GPT de igual profundidad o
algo por debajo, dentro de la desviación entre semillas (nivel 4: 16.5 frente a
16.9\,\% con $N{=}L{=}2$, y 31.6 frente a 34.9\,\% con $N{=}L{=}4$), con 1.8 y 3.4
veces menos parámetros. Con $N{=}8$ iguala o supera al GPT de 8 capas, que tiene
6.7 veces más (nivel 3: 57.2 frente a 54.3\,\%; nivel 4: 45.8 frente a
40.2\,\%). A igual profundidad el tiempo de entrenamiento es el mismo (2.4 frente
a 2.5~h con $D{=}8$): el bucle ahorra parámetros, no tiempo. En C3, el Looped
$N{=}4$ supera al GPT con $d{=}128$ en los niveles 3 y 4 (47.0 frente a 35.0\,\% y
31.6 frente a 24.3\,\%), con mucha menos varianza entre semillas.

\paragraph{Generalización a más operaciones.}
En los niveles 5 y 6 la ventaja del Looped $N{=}8$ es mayor: 33.4 y 21.2\,\%
frente a 25.4 y 13.6\,\% del GPT de 8 capas, que apenas mejora al de 4 capas
(23.9 y 14.5\,\%) aunque sí lo supera en los niveles 3 y 4. Una posible causa es
que, con como máximo 4 operaciones en el entrenamiento, nada obliga a las capas
adicionales del GPT a repetir un mismo paso: ganan capacidad para ajustar los
niveles 1 a 4, mientras que el bucle debe aprender un paso reutilizable. Los
problemas OOD no son más largos que los de entrenamiento (411 caracteres como
máximo en ambos), así que no se trata de posiciones no vistas.

\paragraph{Número de iteraciones necesarias (P2).}
La Fig.~\ref{fig:sweep} varía las iteraciones en test sin reentrenar. Con menos
de las entrenadas la exactitud cae con rapidez, y antes cuanto más difícil es el
nivel: con $N{=}8$, el nivel 1 alcanza 89.6\,\% con 5 iteraciones y 97.1\,\% con 7,
mientras que el nivel 4 sigue mejorando hasta la octava (37.1\,\% con 7 y
45.8\,\% con 8). Las iteraciones necesarias crecen con la dificultad. Con
más de las entrenadas la exactitud también baja, y más deprisa cuanto menor es
$N$ ($N{=}4$ con 8 iteraciones: nivel 1 de 97.3 a 77.7\,\%), y en los niveles 5 y 6 el máximo está en $N$ o en $N{+}1$, por
lo que no observamos extrapolación en profundidad.

\section{Análisis}
\label{sec:analisis}

\paragraph{Logit lens (P2).}
Con \emph{teacher forcing} y 200 problemas por nivel ID (100 OOD), proyectamos cada estado de la posición de inicio de la
respuesta con la cabeza de salida (\emph{logit lens}, \citealp{nostalgebraist2020logitlens}),
$\hat y_t = \arg\max_v[\mathrm{head}(\hid{t})]_v$, y llamamos paso de resolución
al primer $t$ desde el que la respuesta es correcta hasta el final. En el Looped
$N{=}8$ la respuesta se forma de manera gradual y más tarde cuanto mayor es el
nivel: el paso medio de resolución pasa de 4.2 en el nivel 1 a 6.5 en el 4 y a
7.0 en el 6 (Fig.~\ref{fig:interno}, arriba a la izquierda). En el GPT de 8 capas,
en cambio, la respuesta casi no aparece en las capas intermedias (27\,\% con 7
capas en el nivel 1) y surge en la última, sea cual sea el nivel (paso medio 7.6
a 7.8).

\paragraph{Representaciones y atención (P4).}
En ambos modelos el rango efectivo \citep{roy2007effective}, que mide cuántas direcciones usa el estado, baja de unos 145 a unos 46. En el Looped el
cambio relativo $\Delta_t$ se reduce de 4.1 a 0.2, pero la norma crece de 2.5 a
61.5, así que no hay evidencia de un punto fijo. Su atención es casi la misma en
cada iteración (divergencia JS de 0.23 entre las dos primeras y de 0.01 o menos
desde la sexta), mientras que cada
capa del GPT atiende de forma distinta (JS de 0.24 a 0.60). La atención de la posición de respuesta
a los dígitos del enunciado decrece en el Looped cuando la respuesta ya está
formada (nivel 1: de 0.43 a 0.12) y menos en los niveles difíciles (nivel 6: de
0.31 a 0.25).

\paragraph{Tipos de problema y plantillas nuevas.}
En los modelos de profundidad 4 y 8, los problemas de tipo~B (comparación entre
varias personas) con frases desordenadas se resuelven en torno al 40\,\%, frente
al 75--86\,\% de los de tipo~A (relatos de un protagonista o un recipiente). Con redacciones nuevas la exactitud
cae de 64--70\,\% a 23--30\,\%:
los modelos dependen de las plantillas vistas.

\section{Conclusiones y limitaciones}

Un único bloque iterado ocho veces (0.97\,M de parámetros) iguala o supera a un
GPT de ocho capas (6.5\,M) en los niveles difíciles y, sobre todo, fuera de la
distribución, y con los mismos parámetros mejora con claridad al GPT de una capa
(P1, P3). Con dos y cuatro iteraciones queda al nivel del GPT de igual
profundidad o algo por debajo. Las iteraciones necesarias crecen con la
dificultad (P2) y, con menos o con más de las entrenadas, la exactitud cae: no
hay extrapolación en profundidad. Por dentro, el Looped forma la respuesta de
forma gradual, más tarde cuanto mayor es el nivel, con casi la misma atención en
cada iteración, mientras que el GPT la forma en la última capa (P4).

\paragraph{Limitaciones.}
Modelos pequeños y datos sintéticos, que pueden no trasladarse a modelos
grandes ni a texto natural. Los niveles 1 y 2
miden menos razonamiento nuevo: el 100\,\% de sus trazas de test ya aparece en
el entrenamiento con otro texto. Hiperparámetros iguales para todos los modelos, sin ajuste por profundidad.


\bibliography{references}

\end{document}
